\documentclass[12pt,letterpaper]{report}
\usepackage{graphicx}
\usepackage{scrextend}
\usepackage{vmargin}
\usepackage{graphicx}
\usepackage{multirow}
\usepackage[utf8]{inputenc}
\usepackage[spanish]{babel}
\usepackage{multicol}
\usepackage{enumerate}
\usepackage{float}
\usepackage{amsmath, amsthm, amssymb, amsfonts}
\usepackage[usenames]{color}
\parindent=0mm
\pagestyle{empty}
\definecolor{miorange}{rgb}{0.91, 0.43, 0.0}
\begin{document}
\setmargins{2.5cm}      
{1.5cm}                     
{2cm}  
{24cm}                    
{10pt}                          
{1cm}                          
{0pt}                             
{2cm}
\begin{flushright}
\today
\end{flushright}
\vspace{-1.75cm}
\section*{Tarea Clase 11}
\subsection*{Obtenga las soluciones de las siguientes ecuaciones:}
\begin{enumerate}
\item Se quiere colocar un cable desde la cima de una torre de 25 metros altura hasta un punto situado a 50 metros de la base la torre. ¿Cuánto debe medir el cable?
\item Una parcela de terreno cuadrado dispone de un camino de longitud $2\sqrt{2}$ kilómetros (segmento discontinuo) que la atraviesa según se muestra en la siguiente imagen. Calcular el área total de la parcela.
\item La hipotenusa de un triángulo rectángulo mide 10 metros y sus catetos miden x y x+2. ¿Cuánto miden los catetos?.
\item Se desea pintar una cuadrado inscrito en una circunferencia de radio 
$R=3$ como se muestra en la figura. Calcular el área del cuadrado.
\item Calcular cuánto mide el cateto b de un triángulo rectángulo si su otro cateto, a, y su hipotenusa, h, miden 
\begin{eqnarray*}
a=\frac{3\sqrt{2}}{2} \\
h=\frac{\sqrt{194}}{6}
\end{eqnarray*}
\item Hallar las medidas de los lados de una vela con forma de triángulo rectángulo si se quiere que tenga un área de 30 metros al cuadrado y que uno de sus catetos mida 5 metros para que se pueda colocar en el mástil.
\item Si el cateto de un triángulo rectángulo mide x y el otro mide el doble, obtener una fórmula para calcular la longitud de la hipotenusa en función del cateto menor, x.Utilizar la fórmula obtenida para calcular la hipotenusa cuando $x=\sqrt{5}$ y $x=2\sqrt{5}.$
\end{enumerate}
\end{document}